\subsubsection{Laiko eilučių duomenų aibė – CMU Keystroke Dynamics}
\label{subsubsec:keystroke_dynamics}

 Aibę \cite{keystroke_dynamics} sudaro klaviatūros paspaudimų laiko žymos, surinktos iš $51$ tipografo (subjekto), kurių kiekvienas $400$ kartų rinko tą pačią slaptažodžio eilutę \texttt{tie5Roanl} (laikomą stipriu $10$ simbolių slaptažodžiu). Duomenys surinkti per $8$ sesijas (po $50$ pakartojimų kiekvienoje), tarp sesijų išlaikant bent vienos dienos pertrauką, kad būtų užfiksuota natūrali kasdieninė rinkimo ritmo variacija.

Iš viso aibę sudaro $20\,400$ įrašų ir $34$ stulpeliai. Pirmieji trys stulpeliai yra metaduomenys: subjektas (angl. \textit{subject}), pvz., \texttt{s008}, sesijos numeris (angl. \textit{sessionIndex}), $1$–$8$ ir pakartojimo numeris sesijoje \textit{rep}, $1$–$50$. Likę $31$ stulpelis – tai laiko požymiai, susiję su slaptažodžio rinkimu:
\begin{itemize}
    \item \textbf{Hold time} (\texttt{H.key}) – klavišo laikymo laikas (nuo paspaudimo iki atleidimo);
    \item \textbf{Keydown-keydown time} (\texttt{DD.key1.key2}) – laikas tarp dviejų gretimų klavišų paspaudimų;
    \item \textbf{Keyup-keydown time} (\texttt{UD.key1.key2}) – laikas nuo pirmojo klavišo atleidimo iki kito klavišo paspaudimo (gali būti neigiamas).
\end{itemize}
Galioja sąryšis $H + UD = DD$. Visos laiko reikšmės pateiktos sekundėmis, jų matavimo tikslumas $\pm 200$ mikrosekundžių (naudotas išorinis etaloninis laikrodis). Trūkstamų reikšmių aibėje nėra. Pavyzdys ir stulpelių aprašymas pateikti \ref{tab:keystroke_pradiniai_duomenys} bei \ref{tab:keystroke_stulpeliu_aprasymas} lentelėse.

\begin{table}[H]
    \centering
    \caption{CMU Keystroke duomenų aibės struktūra (vienas įrašas, pateikti pirmieji laiko požymiai).}
    \resizebox{\textwidth}{!}{%
    \begin{tabular}{|c|c|c|c|c|c|c|c|}
        \hline
        subject & H.period & DD.period.t & UD.period.t & H.t & DD.t.i & UD.t.i & \ldots \\
        \hline
        s008 & $0{,}3977$ & $0{,}0414$ & $0{,}0329$ & $0{,}4272$ & $0{,}0211$ & $0{,}0180$ & \ldots \\ \hline
    \end{tabular}%
    }
    \label{tab:keystroke_pradiniai_duomenys}
\end{table}
\begin{table}[H]
    \centering
    \caption{CMU Keystroke duomenų aibės stulpelių aprašymas.}
    \begin{tabular}{|l|l|c|}
        \hline
        Stulpelis & Aprašymas & Reikšmių tipas / intervalas \\
        \hline
        \textit{subject} & Subjekto identifikatorius & Tekstinis (pvz., \texttt{s008}) \\ \hline
        \texttt{H.\textit{key}} & Klavišo laikymo trukmė & Realusis sk., $\geq 0$ (s) \\ \hline
        \texttt{DD.\textit{key1}.\textit{key2}} & Laikas tarp dviejų paspaudimų & Realusis sk., $\geq 0$ (s) \\ \hline
        \texttt{UD.\textit{key1}.\textit{key2}} & Laikas nuo atleidimo iki paspaudimo & Realusis sk., $\in \mathbb{R}$ (s) \\ \hline
    \end{tabular}
    \label{tab:keystroke_stulpeliu_aprasymas}
\end{table}

Apibendrinta informacija apie laiko eilučių duomenų aibę pateikta \ref{tab:keystroke_suvestine} lentelėje.

\begin{table}[H]
    \centering
    \caption{CMU Keystroke duomenų aibės bendra suvestinė.}
    \begin{tabular}{|l|l|}
        \hline
        Charakteristika & Reikšmė \\
        \hline
        Šaltinis & CMU (Killourhy \& Maxion, DSN-2009) \\ \hline
        Užduoties tipas & Anomalijų aptikimas / autentifikacija \\ \hline
        Subjektų (klasių) skaičius & $51$ \\ \hline
        Slaptažodis & \texttt{.tie5Roanl} ($10$ simbolių) \\ \hline
        Sesijų skaičius vienam subjektui & $8$ \\ \hline
        Pakartojimų skaičius sesijoje & $50$ \\ \hline
        Iš viso įrašų vienam subjektui & $400$ \\ \hline
        Bendras įrašų skaičius & $20\,400$ \\ \hline
        Stulpelių skaičius & $34$ ($3$ metaduomenys + $31$ požymis) \\ \hline
        Požymių tipai & H, DD, UD laiko žymos (sek.) \\ \hline
        Trūkstamų reikšmių & Nėra \\ \hline
        Failo formatas & \texttt{.csv} \\ \hline
    \end{tabular}
    \label{tab:keystroke_suvestine}
\end{table}